\subsection{An arbitrary idempotent reduction at an extreme corner}

\begin{theorem}[Proved]
\label{CP2:A-extreme-idempotent-corner-reduction}
Let $p$ be an extreme source object in a finite category $C$, put
$M=C(p,p)$, and let $e\in M$ satisfy
\[
 e^2=e,\qquad Me=eMe.
\]
There is an actual finite category $D$ with $BD\simeq BC$, obtained
by replacing the corner $M$ by the monoid $eMe$ with identity $e$,
and each outgoing hom-set $C(p,q)$ by $C(p,q)e$.
All hom-sets between other objects remain unchanged.  In particular,
\[
 |\operatorname{Mor}D|\leq
 |\operatorname{Mor}C|-\bigl(|M|-|eMe|\bigr).
\]
If $e\ne1_M$, this decreases the number of arrows strictly.
For an extreme sink, the dual condition $eM=eMe$ gives the analogous
replacement with incoming hom-sets $eC(q,p)$.
\end{theorem}
\begin{proof}
Adjoin one retract object $p_e$ to $C$, splitting $e$, and call the
result $K$.  This is the full subcategory of the idempotent completion
on the old objects and $p_e$.  Its relevant hom-sets are
\[
 K(p,p_e)=eM,\qquad K(p_e,p)=Me,\qquad
 K(p_e,p_e)=eMe,\qquad K(p_e,q)=C(p,q)e.
\]
For every other old object $q$, $K(q,p_e)$ is empty, because $p$ is
an extreme source.  Composition is inherited from $C$, and the
identity of $p_e$ is $e$.

First the inclusion $i:C\hookrightarrow K$ is a homotopy equivalence.
For an old object $z$, the comma category $(z\downarrow i)$ has
initial object $(z,1_z)$.  For $z=p_e$, an object of
$(p_e\downarrow i)$ is an old object $c$ with an arrow
$f:p_e\to c$, satisfying $fe=f$.  In that comma category put
$o=(p,e:p_e\to p)$.  There is a natural transformation from the
constant functor at $o$ to its identity functor: at $(c,f)$ use the
old arrow $f:p\to c$.  This is a comma morphism because $fe=f$.
For a comma morphism $u:(c,f)\to(d,g)$, naturality is precisely
$uf=g$.  Hence this comma category is contractible.  Quillen's
Theorem~A gives $BC\simeq BK$.

Now let $D$ be the full subcategory of $K$ obtained by deleting the
old object $p$, and let $j:D\hookrightarrow K$.  For a retained
object $z$, the comma category $(j\downarrow z)$ has terminal
object $(z,1_z)$.  In $(j\downarrow p)$ no old retained object
can occur, because no such object has a morphism to $p$.  Its objects
are therefore exactly the arrows $g:p_e\to p$, namely $g\in Me$.
The object $e:p_e\to p$ is terminal.  Indeed a comma morphism
from $g$ to $e$ is an element $k\in eMe$ such that $ek=g$.
Since $e$ is the identity of $eMe$, this requires $k=g$, and this
element belongs to $eMe$ by the hypothesis $Me=eMe$.  Thus there
is exactly one such morphism.  Another application of Quillen's
Theorem~A gives $BD\simeq BK\simeq BC$.

The asserted description of $D$ follows directly from the displayed
hom-sets.  Each new outgoing hom-set is a subset of the corresponding
old hom-set.  If $e\ne1_M$, then $1_M\notin eMe$: an equality
$1_M=eme$ would give $e=e(eme)=eme=1_M$.  Thus its corner is also
strictly smaller.  This proves the arrow bound.  Taking opposite
categories proves the sink statement.
\end{proof}

\begin{corollary}[Proved]
\label{CP2:A-extreme-idempotent-eighteen}
An eighteen-arrow category of homotopy type $\mathbb{CP}^2$ has no
source corner containing a nonidentity idempotent $e$ with $Me=eMe$,
and no sink corner containing one with $eM=eMe$.
In particular neither extreme corner can contain a proper central
idempotent.
\end{corollary}
\begin{proof}
The theorem would produce an actual category with at most seventeen
arrows and the same homotopy type, contradicting
Theorem~\ref{CP2:A-global-lower-eighteen}.  For a central idempotent,
$Me=eM=eMe$, so both versions apply.
\end{proof}

\begin{remark*}
The condition is an exact table predicate: enumerate the nonidentity
idempotents and compare the literal sets $Me$ and $eMe$, or $eM$
and $eMe$ in the sink case.  The earlier minimal-group reduction is
the special case in which $e$ lies in the minimal ideal and $eMe$
is a group.  The present statement does not require minimality or a
group.  It applies to source $3$ in the positive order-six catalog,
where the extra element $a$ is central and idempotent, even though
the minimal ideal has two rows and two columns.
\end{remark*}
